\subsection{Prácticas de Ciclo de Vida de Desarrollo y Aseguramiento de Calidad}\label{subsec:sdlc}

Para garantizar la estabilidad, la mantenibilidad y la evolución segura del ecosistema Paralegales, se implementó un conjunto riguroso de prácticas de ingeniería de software a lo largo del ciclo de vida de desarrollo (\textit{Software Development Life Cycle}, SDLC). Estas prácticas abarcan desde la estandarización del código en los entornos locales de los desarrolladores hasta la validación automática en los repositorios compartidos, conformando un sistema de aseguramiento de calidad continuo y multicapa.

\subsubsection{Estandarización de Código y \textit{Hooks} de Pre-commit}

Con el objetivo de mantener una base de código coherente y libre de errores sintácticos o de estilo, se integraron herramientas de análisis estático y formateo. El ecosistema utiliza ESLint para el \textit{linting} de archivos JavaScript, TypeScript y componentes Vue, y Prettier para el formateo estandarizado y uniforme de múltiples extensiones (incluyendo CSS, JSON, YAML y Markdown).

Para automatizar y hacer obligatorio el cumplimiento de estas reglas antes de que el código ingrese al control de versiones, se adoptó Lefthook como gestor de \textit{hooks} de Git. Su configuración define tres flujos de interceptación clave:

\begin{itemize}
  \item \textbf{\textit{Pre-commit}:} ejecuta en paralelo el formateo (\texttt{prettier --write}) y la validación estática (\texttt{eslint}) exclusivamente sobre los archivos que están en el área de preparación (\textit{staged files}). Si el código es formateado automáticamente, Lefthook lo re-agrega al \textit{commit} sin interrumpir el flujo de trabajo del desarrollador.

  \item \textbf{\textit{Commit-msg}:} garantiza la trazabilidad e historial semántico del proyecto aplicando el estándar \textit{Conventional Commits}. Mediante la integración de Commitlint, cada mensaje de \textit{commit} es evaluado para confirmar que cuenta con una estructura válida (tipo de cambio, alcance y descripción) y que respeta las longitudes máximas configuradas para garantizar legibilidad.

  \item \textbf{\textit{Post-merge}:} detecta si hubo cambios estructurales en el árbol de dependencias (modificaciones en \texttt{pnpm-lock.yaml} o \texttt{package.json}) tras una sincronización o fusión de ramas. Si es el caso, dispara automáticamente la instalación de paquetes (\texttt{pnpm install}) para mantener el entorno local perfectamente sincronizado.
\end{itemize}

\subsubsection{Integración Continua}

Las pipelines de integración continua (\textit{Continuous Integration}, CI) están orquestadas de manera declarativa a través de GitHub Actions. El flujo principal de integración (\texttt{ci.yml}) está diseñado para ejecutarse de forma obligatoria frente a (\textit{Pull Requests}) dirigidas a ramas protegidas, realizando validaciones estrictas con tiempos de procesamiento optimizados:

\begin{itemize}
  \item \textbf{Configuración de entorno y caché:} el \textit{pipeline} aprovisiona contenedores con Node.js e instala eficientemente las dependencias utilizando \texttt{pnpm}, apoyándose en el sistema de caché nativo de GitHub para reutilizar los módulos descargados en ejecuciones previas y reducir la duración total del proceso.

  \item \textbf{Validación diferencial:} se utilizan acciones especializadas de detección de cambios para ejecutar el \textit{linting} y las pruebas de formato únicamente sobre el subconjunto de archivos que sufrieron modificaciones reales, evitando análisis innecesarios sobre código sin cambios.
\end{itemize}

\subsubsection{Estrategia de Pruebas Automatizadas}

La confiabilidad funcional y el blindaje de los datos se soportan en una estrategia de pruebas centrada en la validación de las reglas de seguridad, apoyada en la Firebase Emulator Suite para ejecutar las pruebas en un entorno local y efímero sin afectar entornos reales. Esta estrategia, junto con la validación manual sistemática sobre el entorno de \textit{staging} y la evaluación de pruebas \textit{end-to-end}, se aborda de forma dedicada en el capítulo de pruebas del sistema, por lo que aquí se la referencia sin reproducir su detalle.

La integración de estas prácticas a lo largo del SDLC —desde el \textit{commit} inicial hasta la ejecución del \textit{pipeline} de CI— configura un sistema de aseguramiento de calidad que opera de forma continua y automatizada, reduciendo el riesgo de regresiones y garantizando que cada iteración del sistema mantenga los estándares de calidad y seguridad establecidos para la plataforma.
